\section{System Design}

The system has four layers, shown in Figure 1: (1) a MuJoCo physics simulation of the rescue arena; (2) an MCP server that exposes robot-control tools and the new SAI override tools; (3) a GPT-4o planning agent that calls tools in an agentic loop; and (4) a human supervision interface through which the operator monitors the mission and issues structured overrides.

\subsection{MuJoCo Simulation}
The rescue arena is a bounded 2-D planar world simulated at dt = 0.002 s (500 Hz) with real-time factor 1.0, ensuring that LLM inference latency has genuine operational consequences (delayed assignments translate directly into missed deadlines). The arena contains static obstacles, a configurable number of differential-drive robots, and rescue jobs that arrive stochastically. Each job has a type (medical, structural, fire, evacuation, or hazmat), a location, a priority level (1–5), and a hard deadline in simulation seconds. Robot state—pose, velocity, battery, and current assignment—is serialised into a JSON snapshot at each decision epoch and made available to the MCP server.

\subsection{MCP Server}
The MCP server exposes two categories of tools. The first category—robot-control tools—covers the core planning operations: assign_robot_to_job, query_robot_status, dispatch_move_command, report_job_complete, and get_pending_jobs. These tools were present in our prior system. The second category—the SAI override tools—is introduced in this paper and is described in detail in Section 4. All tool calls pass through a logging middleware that records timestamps, payloads, and round-trip times, providing the raw data for our latency and error-recurrence analyses.

\subsection{Human Supervision Interface}
The supervision interface renders a live top-down view of the MuJoCo arena, a feed of recent LLM tool calls and their outcomes, and an override panel. In the baseline condition (Unstructured), the override panel contains a free-text input box; the operator types a natural-language correction which is appended to the conversation history as a user-turn message. In the SAI condition (Structured), the panel presents a three-tab form—one tab per override type—with labelled fields corresponding to the MCP tool parameters. The operator selects the appropriate tab, fills in the fields, and submits; the interface calls the corresponding MCP override tool. Both conditions are identical in every other respect, including the underlying LLM, simulation, and robot-control tools.

\subsection{Structured Authority Injection}

SAI has three components: an override taxonomy that classifies human corrections, MCP tool schemas for each class, and a persistence model that governs how each class updates the LLM's planning context.